\documentclass{book}
\usepackage[margin=1in]{geometry}
\usepackage{hyperref}
\usepackage{ulem}
\usepackage{tipa} % Add this line to support \textipa command
\hypersetup{
    colorlinks=true,
    linkcolor=black,
    filecolor=magenta,      
    urlcolor=cyan,
    }

\urlstyle{same}

\title{Jian Ming Note}
\author{Junzhang Luo}

\begin{document}

\tableofcontents

\chapter{Leçon 2}
    
    \section{E4}
        \begin{enumerate}
            \item C'est Pascal? Que fait-il?
            \item Est-ce que c'est Nathalie? Que fait-elle?
            \item Qui est-ce? Où est-il?
            \item Qui est-ce? Où est-elle?
        \end{enumerate}

    \section{E5}
        \begin{enumerate}
            \item C'est name.
            \item Oui, c'est name.
            \item Que fait-il?
            \item Il est facteur.
            \item Lina est chercheur.
            \item Pascal est styliste.
        \end{enumerate}

\chapter{Leçon 3}
    \section{E4}
        \begin{enumerate}
            \item Oui, c'est Philippe, Il est facteur.
            \item Oui, il est pilote. Il habite à Berne.
            \item Elle est Sabine. Oui, elle habite à Grenoble.
            \item Oui, il est journaliste. Il habite à Nice.
        \end{enumerate}

    \section{E6}
        \begin{enumerate}
            \item Thomas est facteur. Il habite à Paris.
            \item Victor est journaliste. Elle habite à Vice.
            \item Alice est actrice. Elle habite à Beijing.
            \item Phil est chercheur. Il habite à Bern.
        \end{enumerate}

\chapter{Leçon 4}
    \section{Volcab}
        \begin{itemize}
            \item mouche
            \item soulever
            \item soupirer
            \item la colère
            \item la coque
            \item le col
        \end{itemize}
        
    \section{E4}
        \begin{enumerate}
            \item Où est Philippe habite?
            \item Est-ce que Philippe habite à Nice?
            \item Que fait Marie?
            \item Est-ce que Monsieur Li est avocat?
            \item Est-ce que Li Hong est votre chef de class?
        \end{enumerate}

    \section{E5}
        \begin{enumerate}
            \item Oui, c'est Monsieur Xia.
            \item Oui, il habite à Beijing.
            \item Fanny est actrice.
            \item Oui, Marie est secrétaire.
            \item Pierre est chimiste.
        \end{enumerate}

\chapter{Leçon 5}
    \section{Volcab}
        \begin{enumerate}
            \item un drapeau
            \item une revue
            \item paume
            \item un veau
        \end{enumerate}

    \section{E4}
        \begin{enumerate}
            \item le stylo de Marie
            \item la voiture de Charles
            \item la revue de Pierre 
            \item une robe jaune 
            \item un drapeau rouge
            \item une voiture noir
        \end{enumerate}
        
    \section{E5}
        \begin{enumerate}
            \item Qu'est-ce que c'est?
            \item Qu'est-ce qui c'est?
            \item Est-ce que c'est Philippe?
            \item Est-ce que c'est la jupe de Marie?
            \item Est-ce que René habite à Paris?
        \end{enumerate}
    
    \section{E6}
        \begin{enumerate}
            \item Qu'est-ce que c'est?
            \item C'est une voiture.
            \item Est-ce que c'est la voiture de Marie?
            \item Oui, c'est la voiture de Marie.
            \item Qu'est-ce qui c'est?
            \item C'est Monique. Elle est chimiste.
        \end{enumerate}

\chapter{Leçon 6}   
    \section{Volcab}
        \begin{enumerate}
            \item zèle
            \item épargner
        \end{enumerate}

    \section{Grammar}
        \begin{itemize}
            \item For negation with être, don't change to `de'.
        \end{itemize}

    \section{E5}
        un vélo, un photo, une voiture, un radio, une robe, un stylo, une lettre, un banc, les motos, les manteaux, un roman, les documents

    \section{E6}
        \begin{itemize}
            \item C'est une chaise. C'est la chaise de Sophie.
            \item Est-ce que ce sont des magnétophones?
            \item Ce ne sont pas des magnétophones. Ce sont des radios.
            \item Ce n'est pas le stylo de René.
            \item Ce n'est pas un lac. C'est une rivière.
            \item Ce sont des manteaux de Marie et de Monique.
        \end{itemize}

    \section{E7}
        \begin{enumerate}
            \item Qu'est-ce que c'est?
            \item Ce sont les bancs.
            \item Est-ce que ce sont les magnétophones?
            \item Non, ce ne sont pas les magnétophones, ce sont les radios.
            \item C'est la magnétique de Alice?
            \item Non, ce n'est pas la magnétique de Alice. C'est la magnétique de Pascal.
        \end{enumerate}


\chapter{Leçon 31}

    \section{Phrases}
        \begin{itemize}
            \item à la sortie de (at the end of something)
            \item reconnaître (recognize sth or sbd)
            \item s'empêcher de (to stop oneself from / to refrain from)
        \end{itemize}

    \section{Exercices}
        \begin{enumerate}
            \item Les jeunes font de l'auto-stop quand ils n'ont pas d'argent pour voyager.
            \item Parce qu'on ne connaît pas les gens qu'on fait monter avec nous.
            \item Oui, Albert prenait d'habitude un ou deux auto-stoppeurs quand il allait dans le Midi.
            \item Il essaye de leur parler.
            \item Quand il est à la sortie de Lyon.
            \item Les religieuses avaient une grosse voix.
            \item Les bandits ont attaqué une banque.
            \item Il s'est arrêté devant un hotel et leur a expliqué qu'il n'allait pas plus loin.
            \item Ils ont arrêté une voiture de la police et un gendarme les a reconnus.
        \end{enumerate}
    \section{E2}
        \begin{enumerate}
            \item elle était, nous étions, ils étaient, tu étais, j'étais, vous étiez.
            \item j'ai été, vous avez été, ils ont été, nous avons été, elle a été, tu as été
            \item j'aurai, tu auras, il aura, nous aurons, vous aurez, ils auront
            \item nous avions eu, ils avaient eu, j'avais eu, elle avait eu.
            \item je suis parti / partie, vous êtes parti(s) / partie(s), nous sommes partis / parties, elle est partie
            \item elle était partie, tu étais parti / partie, elles étaient parties, nous étions partis / parties
        \end{enumerate}
    \section{E}
        \begin{enumerate}
            \item Il n'a pas fait ces exercices, parce que il n'a rien compris.
            \item Hier, elle m'a demandé si la police avait déjà arrêté le voleur.
            \item Ce jour-là, il m'a demandé quand j'avais vu Monsieur Dupont.
            \item Quand les enfants n'étaient pas à la maison, elle trouvait le temps long.
            \item Après dîner, l'étudiants ses promanadent en le campus, en les parlent avec le français. 
                    Ensuite, ils retournent leurs les salles de class propres.
                    Certains d'entre eux faisent de devoirs, et l'autures révisent les cours dont leur apprendé avant.
                    Tout les modes sont occupés, ils ne ses couchent pas avant de vingt-trois heures.
            \item Après le dîner, les étuidiants se promenaient sur le campus et parlaient en français.
                    Ensuite, ils retournaient dans leurs salles de classe.
                    Certains faisaient leurs devoirs, et les autres révisaient les leçons qu'ils avaient apprises avant.
                    Tout le monde était occupé, et ils ne se couchaient pas avant vingt-trois heures.
            \item À ce moment-là, les étudiants se  promenaient sur le campus et parlaient en français.
                    Ensuite, ils retournaient dans leurs salles de classe. 
                    Certains commençaient à faire leurs devoirs, et les autres révisaient les leçons qu'ils avaient apprises avant.
                    Tout le monde était occupé, et personne ne se couchait avant onze heures.
            \item Plusieurs étudiants internationnelles voulaient comme les jeunes de françaises, faisaient de voyager par l'auto-stop.
                    Ils attendaient au bord de la route pendant longtemps, mais il n'y a pas des voitures qui avait arrêté pour leur à bord.
                    Ils ses averaient qu'ils ne saivaient pas comme à communicaient à le conducteur.
            \item Quelques étudiants étrangers voulaient voyager en auto-stop comme les jeunes Français.
                    Ils ont attendu longtemps au board de la route, mais aucune voiture ne s'est arrêtée pour les prendre. 
                    En fait, ils ne savaient pas comment faire signe au conducteur.
                \begin{itemize}
                    \item s'arrêter (to stop)
                    \item voyager (travel)
                    \item aucune + noun + ne \dots (no ``pas'' needed, aucune already negative)
                \end{itemize}
        \end{enumerate}

\chapter{Leçon 32}    

    \section{Volcab}
        \begin{itemize}
            \item quant à (as of, about)
            \item en être (to be at a certain stage / a certain state of progress)
            \item où en est-on (where are we? how far along are we?)
            \item se rendre compte de / que (to realize / to become aware of)
            \item prendre qn comme (to take someone as, to employ)
            \item emporter (bring out) $\neq$ apporter (bring in)
            \item propser (à qn) de faire qch
            \item remettre (put back)
            \item remettre qch à qn (give smt to smo)
            \item oser + infinitive (to dare to do something)
        \end{itemize}

    \section{E2}
        \begin{itemize}
            \item être
            \item s'installer
            \item mourir
            \item prendre, sortir
            \item se mettre
            \item croire
            \item naître
            \item vinir
            \item faire
            \item avoir
            \item savoir
            \item être
        \end{itemize}

    \section{E3}
        \begin{itemize}
            \item Il pensait que son fils était paresseux.
            \item Elle croyait qu'elle obtiendrait un poste intéressant.
            \item Nous savions qu'elle n'aimait pas aller au cinéma.
            \item Je étais sûr que vous aviez déjà reçu cette invitation.
            \item Il ne savait pas que vous partiriez demain.
            \item Je pensais qu'il n'y avait pas de cours aujourd'hui.
            \item Il espérait que son ami pourrait venir l'aider ce soir.
            \item Je ne savais pas que ses parents viendraient dans 3 jours.
        \end{itemize}

    \section{E4}
        \begin{itemize}
            \item Il m'a demandé si je partirais le 17 décembre.
            \item Il a demandé à Marie si elle serait à Paris au mois de juin.
            \item Il vous a demandé si vous prendraiez l'avion.
            \item Il a demandé à Pascal s'il aurait un peu de temps libre le lendemain.
            \item Il nous a demandé si nous pouraions réserver une place pour toi.
        \end{itemize}

    \section{E5}
        \begin{itemize}
            \item Le père a demandé à ses fils où ils avaient mis la clé de leur appartement.
            \item La mère a demandé à ses enfants ce qu'ils avaient fait à ce jour-là.
            \item J'ai demandé à Pierre pourquoi Marie avait refusé ces cadeaux.
            \item Le professeur a demandé aux élèves quel film ils avaient vu la veille.
            \item Le directeur m'a demandé quand j'avais appris cette nouvelle.
        \end{itemize}

    \section{E6}
        \begin{itemize}
            \item apportez l'addition
            \item être riche
            \item prendre quelqu'un pour (comme)
            \item rouge de honte
            \item se rendre compte de qc
            \item croire aux explications de qn
        \end{itemize}

    \section{E8}
        \begin{itemize}
            \item Népolen na en 1769 à La Corse, il devenut le emperor de France en 1804, 
                    mais après dix ans, il se retira avec force.
            \item Il m'a dit que la fille prépara toutes ces repas, il n'a fit rien.
            \item Quelqu'un m'a dit que vous êtes en train de ecrit un roman, comme les progrès?
            \item Pendant les cours, l'enseigner a vous informa utiliser le français pour poser un question, 
                    mais pour nous, ce fut très difficile.
            \item Un jour, Albe et sa femme alla une magazine ensemble. Sa femme voit une robe joullie,
                    donc décida en achêter, mais il y avait seulment 30 euros dans la portefeuille, mais le prix de la rope est 35 euros.
                    Quant à Albe, il n'apporta rien d'argent quand il sortit avec sa femme.
        \end{itemize}
        Correction:
        \begin{itemize}
            \item Napoléon naquit en 1769 en Corse et devint empereur de France en 1804, mais dix ans plus tard, il fut forcé de se retirer.
            \item Il m'a dit que ces plats avaient été préparés par sa fille et qu'il n'avait rien fait.
            \item Quelqu'un m'a dit que vous étiez en train d'écrire un roman. Où en êtes-vous?
            \item En classe, le professeur nous a conseillé de poser des questions en français, mais pour nous, c'était assez difficile.
            \item Un jour, Albert et sa femme allèrent ensemble dans un magasin. Sa femme aperçut une jolie jupe et décida de l'acheter.
                    Mais elle n'avait que 30 euors dans son portefeuille, tandis que la jupe coûtait 35 euros. 
                    Quant à Albert, lorsqu'il sortait avec sa femme, il n'emportait jamais d'argent avec lui.
        \end{itemize}

\chapter{Leçon 33}
    \begin{itemize}
        \item rendre (return \& make someone...)
        \item se rendre (surrender)
        \item se rendre à (go to)
        \item avoir l'air (the .adj following is usually property or number)
        \item s'apercevoir de qch (realize)
        \item reporter qch. sur qn (to turn to)
        \item d'une façon générale (generally speaking)
        \item accorder de l'importance à % 對 。。。重視
        \item même after a noun (means itself)
    \end{itemize}

    \section{E2}
        \begin{itemize}
            \item Oui, c'est la salle dont ils ont besoin.
            \item Oui, c'est la personne dont j'ai peur.
            \item Oui, ce sont les fleurs dont elle a envie.
            \item Oui, c'est le visite dont vous parlez.
            \item Oui, ce sont les disques vous avez besoin.
        \end{itemize}

    \section{E3}
        \begin{itemize}
            \item C'est un livre intéressant dont on parle beaucoup.
            \item Ses parents vivent en Italie dont française.
            \item Je suis content dont les nouvelles j'ai reçu.
            \item C'est un grand chanteur dont sa voix est excellente.
            \item C'est un tableau dont il était fier.
            \item C'est Notre-Dame de Paris dont vous avez déjà vu des photos de ce monument.
            \item Pierre et Jacques connaissent bien la Chine dont ils ont visité beaucoup.
            \item Nous avons passé des vacances en montagnes dont je garderai longtemps le souvenir
        \end{itemize}
        Correction
        \begin{itemize}
            \item Bien.
            \item C'est une Française dont les parents vivent en Italie.
            \item J'ai reçu des nouvelles dont je suis content.
            \item C'est un grand chanteur dont la voix est excellente.
            \item C'est un tableau dont le peintre était fier.
            \item C'est Notre-Dame de Paris dont vous avez déjà vu des photos.
            \item Pierre et Jacques connaissent bien la Chine dont ils ont visité beaucoup de villes.
            \item Nous avons passé des vacances en montagnes dont je garderai longtemps le souvenir.
        \end{itemize}
    
    \section{E4}
        \begin{itemize}
            \item C'est une étudiante qui travaille avec moi.
            \item Il n'aime pas la voiture que tu as choisie.
            \item C'est un quartier où il y a beaucoup de magasins.
            \item Ils n'ont pas les livres dont j'ai besoin.
            \item La maison que Paul a fait construire l'an dernier est très belle.
            \item C'est un musicien dont j'oublie toujours le nom.
            \item C'est une région où ils ont baeucoup d'amis.
            \item C'est la personne dont je vous ai parlé hier.
            \item L'automne est une saison où il pleut beaucoup.
            \item Ne prenez pas le sac qui est sur la table.
            \item La robe que Nicole a achetée côute cher.
            \item C'est un endroit où il arrive souvent des accidents.
        \end{itemize}

    \section{E5}
        \begin{itemize}
            \item Oui, je le suis.
            \item Oui, il l'est.
            \item Oui, nous le savons.
            \item Non je ne le sais pas.
            \item Non, nous ne le sommes pas.
            \item Oui, il le désire.
            \item Non, nous ne le savons pas.
            \item Oui, je le pense.
        \end{itemize}

    \section{E6}
        \begin{enumerate}
            \item dont: être satisfait de
            \item où, not qu'
            \item dont les, not dont leurs
            \item l'est, not le est
            \item que, not qui
            \item dont, not qu'
        \end{enumerate}

    \section{E7}
        \begin{enumerate}
            \item À ce moment là, elle a l'air en bonne santé, mais pas en réalité.
            \item Je peux vous retourne l'ordinateur dans le après-midi.
            \item Le réunion est reporté à la semaine prochaine.
            \item Ne pas de raison, il est tard toujours.
            \item C'est le scénario bizarre dont tout le monde parle.
            \item Qu'est-ce que vous me direz quoi dont vous savez besoins?
        \end{enumerate}
        Corrected:
        \begin{enumerate}
            \item À ce moment-là, elle avait l'air en bonne santé, mais en réalité, ce n'était pas le cas.
            \item Je peux vous rendre l'ordinateur cet après-midi.
            \item La réunion est reportée à la semaine prochaine.
            \item On ne sait pas pour quelle raison, il est souvent en retard.
            \item C'est un phénomène étrange dont tout le monde parle actuellement.
            \item Pouvez-vous me dire ce dont vous avez besoin.
        \end{enumerate}
        ce \dots dont (what \dots need)

\chapter{Leçon 34}
    \begin{itemize}
        \item avec (can be used to express reason)
        \item ne \dots plus que (only)
        \item travailler à (devoted to)
        \item devoir + infinitive (to be expected to)
        \item en peser + number (to weigh that number)
        \item use ``en'' to replace a quantity noun, or something already mentioned after a number.
        \item renseigner qn (tell qn)
        \item se renseigner % 打聽情況
        \item couvrir % 覆盖，盖过，完成（路程）
        \item servir à qc (to serve for something)
        \item À quoi \dots sert? (What is \dots for)
        \item se mettre à (start doing something)
        \item Pronoun: ``chacun(e)'' to replace ``chaque + noun'', instead of repeating (usually mean each one)
            \begin{itemize}
                \item Chaque + singular noun
                \item Chacun(e) de + determiner (the/these/my/etc.) + plural noun phrase
            \end{itemize}
        \item être de + number / amount
        \item s'enorgueillir (de) (be proud of \dots)
        \item En surtout (most importantly, above all)
    \end{itemize}
    Le participe présent: take -ons and replace with -ant. Usage:
    \begin{enumerate}
        \item same as qui + verb
        \begin{center}
            On recrute des interprètes connaissant à la fois le français et l'anglais.
            \\
            connaissant = qui connaissent
        \end{center}
        \item Check comment % 表示時間或原因的狀語從句
        \begin{center}
            Traversant la rue, j'ai rencontré Paul = Quand je traversais la rue, \dots
        \end{center}
        \item Use as adjective, but agree to gender.
        \begin{center}
            La situation est encourageante. Je suis en train de lire des livres intéressants.
        \end{center}
    \end{enumerate}
    Tout:
    \begin{enumerate}
        \item Modifies a noun. agrees to gender.
        \item Replace noun/group, agrees to gender.
        \item Modifies adj / adv, no gender change, except for feminine adj that starts with consonant and silent h (toute is used).
    \end{enumerate}

    \section{E2}
        \begin{enumerate}
            \item Danielle Fabre est très connue pour ses émissions de radio.
            \item Alain Gautier est très connu pour ses films policiers.
            \item Cette journaliste est très connue pour son livre sur la Chine.
            \item Ces scientifiques sont très connus pour leur voyage au Pôle Sud.
        \end{enumerate}

    \section{E3}
        \begin{itemize}
            \item savoir: sachant
            \item venir: venant
        \end{itemize}

    \section{E4}
        \begin{enumerate}
            \item On forme dans cet institut des étudiants connaissant à la fois le français et l'anglais.
            \item Les élèves ayant beaucoup de devoirs à faire ne peuvent pas se coucher avant 11 heures du soir.
            \item Les personnes travaillant dans ce laboratoire sont toutes très jeunes.
            \item Les amis étrangers chantant la Marseillaise travaillent dans notre entreprise.
            \item J'ai lu un article exposant de grands problèmes internationaux.
        \end{enumerate}

    \section{E5}
        \begin{enumerate}
            \item Ma fille étant malade, elle n'est pas allée à l'école la semaine denière.
            \item Il ayant des lettres à écrire, il ne va pas au cinéma.
            \item Je venant faire mes courses, je rencontre toujours Mme Dupont.
            \item Il ayant peur d'être en retard et il se met à courir.
            \item Elle ne pouvant pas venir à notre soirée, Monique nous a téléphoné pour s'excuser.
        \end{enumerate}
        Correction:
        \begin{enumerate}
            \item Ayant des lettres ... (remove the subject for participe présent)
            \item En allant faire mes courses (quand expresses time, so we use ``en + participe présent'')
            \item Ayant peur d'être
            \item Ne pouvant pas 
        \end{enumerate}
    
    \section{E6}
        \begin{enumerate}
            \item Chaune de ces salles a 150 places.
            \item Chaune de ces usines emploie environ 200 ouvriers.
            \item Chaun de ces moteurs est de 23,000 C.V.
            \item Chaune de ces valises pèse 25 kilos.
        \end{enumerate}

    \section{E7}
        L'Airbus est un grand vol, construitant par des plusieurs entreprises de France, Alemane, Angeleterre et Espagne.
        Il a mesure 53.62 mètres de long et 44.84 mètres de large, pèse 142 tonnes. 
        Il a vitesse chaque heure est 950 kilomètres, il a emporte 345 passagère.
        L'Airbus a vol pour un première fois sur octobre vingt-huit en 1972. 
        Maitenant, il est en concurrence sérieuse vers le Boeing aux Etáts-Unis. 
        \\
        Correction:
        \\
        L'Airbus est un grand avion de ligne, construit par plusieurs entreprises françaises, allemandes, britanniques et espagnoles.
        Il mesure 53,62 mètres de long et 44,84 mètres de large, et pèse 142 tonnes.
        Sa vitesse est de 950 kilomètres à l'heure, et il peut transporter 345 passagers.
        L'Airbus a volé pour la première fois le 28 octobre 1972.
        Maintenant, il est en forte concurrence avec le Boeing fabriqué aux États-Unis.
        \\
        More polished version:
        \\
        L'Airbus est un grand avion de ligne fabriqué par plusieurs entreprises françaises, allemandes, britanniques et espagnoles. 
        Il mesure 53,62 mètres de long et 44,84 mètres de large, et pèse 142 tonnes. 
        Sa vitesse est de 950 kilomètres à l'heure, et il peut transporter 345 passagers. 
        L'Airbus a effectué son premier vol le 28 octobre 1972. 
        Aujourd'hui, il est en forte concurrence avec le Boeing fabriqué aux États-Unis.

    \section{E8}
        \begin{enumerate}
            \item pronoun (replaces a noun/object)
            \item pronoun (refers to ils, means all of them)
            \item pronoun (stand along)
            \item adverb
            \item adj. modifies the noun
            \item pronoun
        \end{enumerate}

    \section{E9}
        \begin{enumerate}
            \item J'aime toutes les langues.
            \item Tous les étudiants travaillent sérieusement.
            \item Est-ce que vous avez compris tout le discours?
            \item Est-ce que tous les professeurs partent à 6 heures?
            \item Je lis des journaux tout le matin.
            \item Nous avons dansé pendant toute la soirée.
        \end{enumerate}

\chapter{Leçon 35}
    \begin{itemize}
        \item en avoir pour % 为此需要
        \item regarder à qch (consider, pay attention to)
        \item se servir de % 利用，使用
        \item tarder à + inf % 延迟， 延缓
        \item à cela (to that / about that)
        \item rendre la joie à qn (bring happiness to someone)
        \item il s'agit de (regarding)
        \item d'après (according to)
        \item avoir mal à (not feeling well for some body parts)
        \item il vaut mieux + inf (it is better to)
        \item il est question de (it is about, it is a question of)
        \item tarder à + inf (delay)
        \item revenir (return)
        \begin{itemize}
            \item revenir à qch (return to)
            \item revenir de qch (recover, get out of)
        \end{itemize}
        \item For ``si'' expressions, use imparfait for unlikely hypothesis, for likely ones, just use present tense.
        \item de quoi (enough)
        \item chaque fois que (everytime, when)
    \end{itemize}

    \section{E2}
        \begin{enumerate}
            \item Je voudrais deux kilos de pommes.
            \item Nous voudrions parler au directeur.
            \item Il revendrait dans trois jours, s'il le pouvait.
            \item Vous réussiriez à ces examens, si vous travailliez plus.
            \item Pourriez-vous m'aider?
            \item Nous ferions de l'auto-stop, si nous voyagions en Europe.
        \end{enumerate}

    \section{E3}
        \begin{enumerate}
            \item Si tout le monde était là, on partirait tout de suite.
            \item Si nous étions à Paris, nous visiterions le palais du Louvre.
            \item Si vos amis venaient vous voir, vous seriez heureuse?
            \item S'il y avait de la neige, nous ferions du ski.
            \item Si vous répétiez ces mots, nous les sauriez.
            \item S'ils écoutaient la radio française, ils comprendraient mieux les actualités en France.
            \item Si nous étions pressés, nous prendrions l'avion.
            \item S'il avait beaucoup d'argent, il voyagerait au Pôle Sud.
            \item Si nous proposions cette solution, les autres ne seraient pas d'accord.
            \item Si on les faisait travailler plus, ils demanderaient une augmentation de salaire.
        \end{enumerate}

    \section{E4}
        \begin{enumerate}
            \item J'apprendrais davantage le français.
            \item Je consulterais un médecin.
            \item Je voyagerais autour du monde.
            \item Je visiterais Paris et le musée du Louvre.
            \item Je refuserais poliment.
        \end{enumerate}
    
    \section{E5}
        \begin{enumerate}
            \item S'il pleuvait, on prendrait un parapluie.
            \item Si le professeur ne vient pas, on devra annuler cet événement.
            \item Si un Martien arrivait, je le saluerais.
            \item Elle arrêtera de travailler, si elle a assez d'argent.
            \item Paul partirait à l'étranger, s'il voulait.
            \item Nous irions au cinéma, si nous avions du temps.
        \end{enumerate}

    \section{E6}
        \begin{enumerate}
            \item use imparfait: voyais
            \item S'il pleut demain, nous resterons à la maison.
            \item Si marie était malade, elle n'irait pas au cinéma.
            \item Si vous avez besoin de quelque chose, vous me téléphonerez.
            \item Je réfléchirais davantage ce problème, si j'avais beaucoup de temps.
        \end{enumerate}

    \section{E7}
        \begin{enumerate}
            \item Nous travaillons en chantant.
            \item La mère répondait en souriant.
            \item Ils racontent leur histoire en mangeant.
            \item Le chef explique le problème en faisant des gestes.
            \item J'ai dit au revoir en sortant.
            \item L'enfant est entré en pleurant.
        \end{enumerate}

    \section{E8}
        \begin{enumerate}
            \item Les étudiants apprennent cette chanson en écoutant un enregistrement.
            \item Je suis sorti de la classe en fermant la porte et les fenêtres.
            \item Elle a appris ces nouvelles en regardant la télévision.
            \item Il a pris froid en marchant sous la pluie.
            \item Ils rentrent à la maison en prenant l'autobus.
            \item Il s'est coupé le doigt en faisant la cuisine.
        \end{enumerate}

    \section{E9}
        \begin{enumerate}
            \item Ils déjeunent en attendant leurs amis.
            \item N'ayant pas assez d'argent, ils n'ont pas pu acheter ce poste de télévision.
            \item En discutant, nous avons pu prendre rapidement une décision.
            \item Quand je suis arrivé, j'ai vu le directeur entrant dans son bureau. (the verb is not for ``je'')
            \item On n'arrivera à l'heure qu'en prenant un taxi. (by taking a taxi)
            \item Habitant loin de son usine, il arrive quelquefois en retard.
            \item En lisant plusieurs romans de cet auteur, on le connaîtra mieux.
            \item Vous obtiendrez de meilleurs résultats, en faisant plus d'efforts.
        \end{enumerate}

    \section{E10}
        \begin{enumerate}
            \item Pierre a un peu mal à la tête, sa mère le conselle trouver à docteur Morland.
            \item Il est sen enfant, il est excité simplement.
            \item Veuillez relire cet article, il prends au plus 10 minutes.
            \item Si je ne la envoyait une message, elle s'inquiète.
            \item Si vous ne regarderez pas ce match de football demain, nous donnerons ces tickets à l'autres étudiants.
            \item Si vous ne me répondiez pas, nous consultosn l'autres avocats.
            \item Ne lisez pas le livre en écoutant la radio.
            \item Le professeur a explique la grammare en donnant quelque phrases exemples.
        \end{enumerate}
        Correction:
        \begin{enumerate}
            \item Pierre a un peu mal à la tête, sa mère lui conseille d'aller voir le docteur Morland.
            \begin{itemize}
                \item ``lui'' because conseiller à Pierre (indirect object)
            \end{itemize}
            \item C'est un enfant sensible, il s'énerve très facilement.
            \item Veuillez relire cet article, cela ne prendra que dix minutes au maximum.
            \item Si je ne lui envoie pas de message, elle sinquiétera.
            \item Si vous ne pouvez pas aller voir ce match de footbal demain, nous donnerons ces billets à d'autres étuidiants.
            \item Ne lisez pas en écoutant la radio.
            \item Le professeur a expliqué la grammaire en donnant quelques phrases d'exemple.
        \end{enumerate}

\chapter{Lecon 36}
    \begin{itemize}
        \item réaliser des ventes (to generate sales)
        \item se réjouir de qch (be pleased for something)
        \item être inférieur à (lower than)
        \item être supérieur à (higher than)
        \item par suite de (since)
        \item vaincu (Defeated, overcome, beaten)
        \item finir par + inf (to end up do something)
        \item le sien (his own), check the chart on page 84.
        \item alors que % 轉折
        \item jeter un coup d'œil % 打量, 对...瞟一眼
        \item à l'avance (in advance)
        \item après tout (regardless)
        \item ça fait \dots que (have been)
        \item être en panne (broken)
        \item avoir lieu (take place)
        \item en faire autant (do the same)
        \item se faire (to be done / to happen)
        \item céder le pas à quelque chose / quelqu'un (to give way to something/someone)
        \item il s'agit de (it is about / it is a matter of)
        \item se rendre compte que (to realize that)
        \item se tendre (to become tense / to tighten)
        \item aussi (therefore, must have inversion to follow up)
    \end{itemize}

    \section{E2}
        \begin{enumerate}
            \item Votre fille est devenue plus intelligente que la mienne.
            \item Ma maison paraît plus ancienne que la vôtre.
            \item Leur appartement semble plus propre que le vôtre.
            \item Notre travail n'est pas plus important que le leur.
            \item Ma place est meilleure que la tienne.
            \item Mon chien a l'air plus calme que le sien.
            \item Votre voiture paraît plus petite que la sienne.
            \item Ta mère est-elle plus âgée que la mienne.
        \end{enumerate}
        
    \section{E3}
        \begin{enumerate}
            \item As-tu trouvé la tienne?
            \item Marie a pris la sienne?
            \item Le boulanger a fermé le sien?
            \item Il parle souvent de la sien?
            \item M. Morin a besoin de la vôtre?
            \item M. Dupont et M. Renou sont content de les leurs?
        \end{enumerate}
        Correction:
        \begin{enumerate}
            \item Good
            \item Marie a-t-elle pris le sien?
            \item Le boulanger a-t-il fermé le sien?
            \item Sa femme parle-t-elle souvant du sien?
            \item Avez-vous besoin de la vôtre?
            \item M. Dupont et M. Renou sont-ils contents des leurs?
        \end{enumerate}

    \section{E4}
        \begin{enumerate}
            \item Je fais mes exercices, tu fais les tiennes?
            \item Je parle de mon travail, tu parles du tien?
            \item Pierre lit son journal, Marie lit le sien?
            \item Nous lisons nos journaux, vous lisez les vôtres?
            \item Je te donne ma adresse, peux-tu me donner la tienne?
            \item Elle pense souvent à ses parents; et vous, est-ce que vous pensez souvent à les vôtres?
        \end{enumerate}
        Correction:
        \begin{enumerate}
            \item les tien, exercice is masculine.
            \item Good
            \item Good
            \item Good
            \item mon adresse, because starts with a 
            \item à les $\to$ aux.
        \end{enumerate}

    \section{E5}
        \begin{enumerate}
            \item Il y a deux livres sur la table, ce sont le mien. Ceux qui sont sur la table, ce sont le mien.
            \item Il y a une voiture devant la maison, c'est la leur. Celle qui est devant la maison, c'est la leur.
            \item Il y a des lettres sur le bureau, ce sont les vôtres. Celles qui sont sur le bureau, ce sont les vôtres.
            \item Il y a un vélo dans la cour, c'est le sien. Celui qui est dans la cour, c'est le sien.
        \end{enumerate}
        Correction:
        \begin{enumerate}
            \item 
            \begin{enumerate}
                \item Il y a deux livres sur la table, ce sont les miens.
                \item Ceux qui sont sur la table, ce sont les miens.
            \end{enumerate}
        \end{enumerate}

    \section{E6}
        \begin{enumerate}
            \item Nous faisons nos exercies, il faites les siens.
            \item Voilà deux billets, celui-ci est à moi, celui-la est à toi.
            \item Les livres de poche dont on trouve partout et ils ne sont pas chers poussent les Français à la lecture.
            \item Ce livre est très intéressant, pouvez-vous me donnez prêter?
            \item Cette jolie robe est à Marie. Sa mère vient de un offrir pour son anniversaire.
        \end{enumerate}
        Correction:
        \begin{enumerate}
            \item Nous faisons nos exercices, vous faites les vôtres.
            \item Voilà deux billets, celui-ci est à moi, celui-là est à toi.
            \item Les livres de poche qu'on trouve partout et qui ne sont pas chers poussent les Français à la lecture.
            \item Ce livre est très intéressant, pouvez-vous me le prêter ?
            \item Cette jolie robe est à Marie. Sa mère vient de la lui offrir pour son anniversaire.
        \end{enumerate}

    \section{E7}
        \begin{enumerate}
            \item Votre fils étudie tellement sérieux, mais le mien ne connait que jouer.
            \item Marie a emporté mon enregistreur, le sien est cassé.
            \item Passé, le père de Charolette est un agriculteur, maitenant Charolette est également un agriculteur,
                    son fils a-t-il devenirait un agriculteur?
            \item Celles légumes est très cher, ne achéte pas d'ici, 
                    parce que à la face de chez moi, on peut les achéte pour la moitié d'argent.
            \item Après l'énemi avait défaiter, il va capituler finalement.
        \end{enumerate}
        Correction:
        \begin{enumerate}
            \item Votre frère travaille très sérieusement, mais le mien ne pense qu'à s'amuser.
            \item Marie a emporté mon magnétophone, car le sien est en panne.
            \item Autrefois, le père de Charles était paysan ; maintenant, Charles est aussi paysan. Son fils le sera-t-il encore à l'avenir ?
            \item Ces légumes sont trop chers. Ne les achetez pas ici, car en face de chez nous, on peut acheter les mêmes pour moitié prix.
            \item Après avoir été vaincu, l'ennemi a fini par capituler.
        \end{enumerate}
    
\chapter{Leçon 31-36 Revision}
    \section{Conjugation}
        \begin{itemize}
            \item parlerai, parlerais, avait parlé
            \item dira, dirions, aviez dit
            \item partira, partirais, étaient partis
            \item finiras, finirait, avions finis 
            \item prendra, prendrons, aviez pris
            \item venirai, venirais, avait voit 
            \item va, allerait, avait alle
            \item oserai, oserais, avait osé
            \item se levera, se leverait, avions les levé
        \end{itemize}
        Correction:
        \begin{itemize}
            \item parler : je parlerai, je parlerais, il avait parlé
            \item dire : il dira, nous dirions, vous aviez dit
            \item partir : elle partira, tu partirais, ils étaient partis
            \item finir : je finirai, il finirait, nous avions fini
            \item prendre : il prendra, nous prendrions, vous aviez pris
            \item venir : tu viendras, tu viendrais, elle était venue
            \item aller : il ira, il irait, elle était allée
            \item oser : tu oseras, j'oserais, il avait osé
            \item se lever : il se lèvera, elle se lèverait, nous nous étions levés
        \end{itemize}

    \section{Remplissez le tableau}
        \begin{itemize}
            \item faire, fait, faisant
            \item venir, vient, viennant
            \item savoir, su, savant
            \item voir, vu, voyant
            \item partir, parti, partant
            \item avoir, eu, ayant
            \item être, dû, étant
            \item pouvoir, peut, pouvant
        \end{itemize}
        Correction:
        \begin{itemize}
            \item faire, fait, faisant
            \item venir, venu, venant
            \item savoir, su, sachant
            \item voir, vu, voyant
            \item partir, parti, partant
            \item avoir, eu, ayant
            \item devoir, dû, devant
            \item pouvoir, pu, pouvant
        \end{itemize}

    \section{Convenables}
        \begin{enumerate}
            \item Qu'est-ce qu'il fait hier?
            \item Si nous étions en France, il nous étions plus facile d'apprendre le français. 
            \item Il m'a demandé si je irais le voir le lendemain.
            \item Il m'a demandé si je voirais son fils la veille.
            \item Il m'a dit qu'il était à la campagne ce jour-là.
            \item Si nous prendrions le Concorde, nous arriverions plus tôt à New York.
            \item Il pensait à la soirée qu'il avait passé la veille.
            \item Elle était inquitète, parce qu'elle ne recevoirait pas de nouvelles de son père.
        \end{enumerate}
        Correction:
        \begin{enumerate}
            \item Qu'est-ce qu'il a fait hier ?
            \item Si nous étions en France, il nous serait plus facile d'apprendre le français.
            \item Il m'a demandé si j'irais le voir le lendemain.
            \item Il m'a demandé si j'avais vu son fils la veille.
            \item Il m'a dit qu'il était à la campagne ce jour-là.
            \item Si nous prenions le Concorde, nous arriverions plus tôt à New York.
            \item Il pensait à la soirée qu'il avait passée la veille.
            \item Elle était inquiète, parce qu'elle n'avait pas reçu de nouvelles de son père.
        \end{enumerate}
        For avoir, if the direct object comes before the verb, then agree to gender, 7 is an example, where ``qu'' is ``la soirèe''

    \section{Pronom Relatif}
        \begin{enumerate}
            \item C'est un garage où se trouve près de chez moi.
            \item Je ne connais pas livre dont vous avez parlé.
            \item Je n'ai pas lu les articles que vous avez écrits.
            \item Je connais bien le médecin qu'il a consulté l'autre jour.
            \item Dans la région que j'habite on cultive du blé.
            \item Dans la région dont je viens la culture principale est le riz.
            \item On cherche la personne qui la voiture s'est arrêtée devant la porte du magasin.
            \item Les plats que votre femme a préparés sont excellents.
            \item La robe qu'elle a essayée est trop longue.
            \item Un livre est un cadeau qui fait toujours plaisir.
        \end{enumerate}
        Corrections:
        \begin{enumerate}
            \item C'est un garage qui se trouve près de chez moi.
            \item Je ne connais pas le livre dont vous avez parlé.
            \item Je n'ai pas lu les articles que vous avez écrits.
            \item Je connais bien le médecin qu'il a consulté l'autre jour.
            \item Dans la région où j'habite, on cultive du blé.
            \item Dans la région d'où je viens, la culture principale est le riz. (venir de quelque part, to come from somewhere)
            \item On cherche la personne dont la voiture s'est arrêtée devant la porte du magasin.
            \item Les plats que votre femme a préparés sont excellents.
            \item La robe qu'elle a essayée est trop longue.
            \item Un livre est un cadeau qui fait toujours plaisir.
        \end{enumerate}

    \section{Transformez}
        \begin{enumerate}
            \item Étant nombreux, on s'amuse mieux.
            \item Étant gentil, on se fait des amis.
            \item Étant imprudent, on cause des accidents.
            \item Étant sportif, on reste jeune.
            \item Étant bavard, on gêne les autres.
        \end{enumerate}

    \section{Transform}
        \begin{enumerate}
            \item N'ayant pas faim, je ne mange pas.
            \item N'ayant pas soif, je ne boire pas.
            \item N'ayant pas des lunettes, je ne peux pas lire.
            \item N'ayant pas le temps, je ne peux pas attendre.
            \item N'ayant pas un billet, je ne peux pas entrer.
        \end{enumerate}
        Correction:
        \begin{enumerate}
            \item N'ayant pas soif, je ne bois pas.
        \end{enumerate}

    \section{Refaites}
        \begin{enumerate}
            \item Si il connaît l'espagnol, il pourait donc vous traduire cet article.
            \item Si j'avais son numéro de téléphone, je pourais donc lui passer un coup de fil.
            \item Si il y avait de neige dans les montagnes, nous pourrions donc faire du ski.
            \item Si vous étiez sportif, vous pourriez donc faire une si longue course.
            \item Si j'avais de temps, je pourais donc vous accompagner à l'aéroport.
            \item Si nous avions un grand appartement, nous pourrions donc vous recevoir tous.
        \end{enumerate}
        Correction:
        \begin{enumerate}
            \item S'il connaissait l'espagnol, il pourrait vous traduire cet article.
            \item Si j'avais son numéro de téléphone, je pourrais lui passer un coup de fil.
            \item S'il y avait de la neige dans les montagnes, nous pourrions faire du ski.
            \item Si vous étiez sportif, vous pourriez faire une si longue course.
            \item Si j'avais le temps, je pourrais vous accompagner à l'aéroport.
            \item Si nous avions un grand appartement, nous pourrions vous recevoir tous.
        \end{enumerate}

    \section{Reliez}
        \begin{enumerate}
            \item C'est un beau jardin dont M. Legrand est propriétaire.
            \item La concierge s'appelle Duroc dont les locataires sont mécontents.
            \item Nous avons vu une pièce de théâtre dont je connais l'auteur.
            \item Nous avons traversé une forêt dont nous ne connaissons pas le nom.
            \item Ce sont des plantes médicinales dont on a parfois besoin.
            \item Il sortit de l'appartement dont il ferma à clé la porte.
        \end{enumerate}
        Correction:
        \begin{enumerate}
            \item C'est un beau jardin dont M. Legrand est propriétaire.
            \item La concierge dont les locataires sont mécontents s'appelle Duroc.
            \item Nous avons vu une pièce de théâtre dont je connais l'auteur.
            \item Nous avons traversé une forêt dont nous ne connaissons pas le nom.
            \item Ce sont des plantes médicinales dont on a parfois besoin.
            \item Il sortit de l'appartement dont il ferma la porte à clé.
        \end{enumerate}

    \section{Choisissez}
        \begin{enumerate}
            \item Je n'y étais pas, parce qu'on ne m'avait pas invité.
            \item Le ministre de l'Education a publié un décret modifiant les dates des vacances scolaires.
            \item J'ai déjà remis mes devoirs au chef de classe. As-tu remis les tiens?
            \item c
            \item a
            \item a
            \item b
            \item c
        \end{enumerate}
        Correction:
        \begin{enumerate}
            \item A. Je n'y étais pas, parce qu'on ne m'avait pas invité.
            \item C. Le ministre de l'Éducation a publié un décret modifiant les dates des vacances scolaires.
            \item B. J'ai déjà remis mes devoirs au chef de classe. As-tu remis les tiens ?
            \item A. Mon frère étant malade, il n'est pas allé à l'école ce jour-là.
            \item C. Il m'a dit qu'il avait déjà su cette nouvelle la veille.
            \item C. Si j'étais riche, je vous prêterais de l'argent.
            \item B. Nathalie, sachant que ses parents rentrent tard, reste dans la classe et y fait ses devoirs.
            \item C. Mon ordinateur ne marche plus. Pouvez-vous me prêter le vôtre ?
        \end{enumerate}

    \section{Volcab}
        \begin{itemize}
            \item se servir de (use)
            \item qu'à cela ne tienne (not a big deal)
        \end{itemize}

\chapter{Lecon 37}
    \begin{itemize}
        \item à la légère % 草率的
        \item il me paraît impossible de (does not have a subject, means ``from my perspective, we shouldn't')
        \item tirer un avantage de qch. (get advantage from)
        \item y compris / non compris (include / not include)
            \begin{itemize}
                \item if placed after the noun, then agree to gender
                \item ex. les voitures non comprises
            \end{itemize}
        \item considérer que (think (opinion))
        \item considérer qch. comme (view someone as)
        \item s'assurer
            \begin{enumerate}
                \item insurance: il s'est assuré contre l'incendie.
                \item obtain: Je voudrais m'assurer d'une bonne place.
                \item investigate / confirm: assurez-vous de l'exactitude de cette nouvelle.
            \end{enumerate}
        \item contraire à (not good for / contrary)
        \item paraître
            \begin{enumerate}
                \item appear: Le soleil commence à paraître. Un avion parut dans le ciel.
                \item seems: Il paraît approuver cette idée. Le voyage paraît très long.
                \item seems (impersonal): Il paraît qu'on va doubler les impôts. Il me paraît préférable que vous sortiez.
            \end{enumerate}
        \item tout d'un coup (sudden)
    \end{itemize}

    \section{Le Subjonctif Présent}
        \begin{enumerate}
            \item Remove -ent from `ils' form, add: -e, -es, -e, -ions, -iez, -ent. For nous and vous form, use their corresponding stems.
            \item Irregular: avoir (aie), être (sois), aller (allie), faire (fasse), pouvoir (puisse), savoir (sache), vouloir (veuille), falloir (faille), pleuvoir (pleuve).
        \end{enumerate}

    \section{E2}
        \begin{enumerate}
            \item Il faut que vous fassiez ces exercices.
            \item Il est utile qu'ils apportent cette carte.
            \item Il est temps que tu ailles à l'école.
            \item Il est nécessaire que vous prendiez ces médicaments.
            \item Il est important que nous connaîtions ces nouveaux mots.
            \item Je doute qu'il revenisse.
            \item Je veux que vous réfléchissiez avant de répondre à ces questions.
            \item Ils regrettent que nous ne puissent pas y aller avec eux.
            \item Je suis désolé qu'ils n'aient pas assez de temps.
            \item Je suis heureux que vous soyez en bonne santé.
        \end{enumerate}
        vous preniez, nous connaissions, il revienne, nous ne puissions pas

    \section{E3}
        \begin{enumerate}
            \item Je désire que vous ditiez bonjour à tous ceux que vous rencontrerez dans ce bâtiment.
            \item Le chef exige que nous soyions à l'heure.
            \item Le professeur voudrait que les étuidiants sachent employer le subjonctif.
            \item Ses parents souhaitent qu'elle obtienne ce diplômee.
            \item Nous aimerions que MArie fasse un voyage en Chine.
        \end{enumerate}
        Correction:
        \begin{enumerate}
            \item Je désire que vous disiez bonjour à tous ceux que vous rencontrerez dans ce bâtiment.
            \item Le chef exige que nous soyons à l'heure.
            \item Le professeur voudrait que les étudiants sachent employer le subjonctif.
            \item Ses parents souhaitent qu'elle obtienne ce diplôme.
            \item Nous aimerions que Marie fasse un voyage en Chine.
        \end{enumerate}

    \section{E4}
        \begin{enumerate}
            \item Pierre a peur qu'il pleuve.
            \item Le directeur craint qu'ils ne finissent pas ce travail avant 6 heures.
            \item Nous sommes contents qu'ils aient de longues vacances.
            \item Nous regrettons que les exercices soient trop difficiles.
            \item Les touristes étrangers sont surpris que les magasins soient fermés à 18 heures.
        \end{enumerate}

    \section{E5}
        \begin{enumerate}
            \item Il est possible que nous vienons demain.
            \item Il est indispensable que vous apportiez vos diplômes.
            \item Il est regrettable que vous ne parliez pas anglais.
            \item Il vaudrait mieux que vous aillez le voir tout de suite.
            \item Il est normal que les touristes posent des questions au guide.
        \end{enumerate}
        Corrections (1,4):
        \begin{enumerate}
            \item Il est possible que nous venions demain.
            \item Il vaudrait mieux que vous alliez le voir tout de suite.
        \end{enumerate}

    \section{E6}
        \begin{enumerate}
            \item Je suis content que je prenne des vacances.
            \item Il est content que vous veniez demain.
            \item Je suis contente qu'elle soie d'accord avec vous.
            \item Elle a peur qu'elle sortie seule.
            \item Nous sommes heureux que vous puissiez nous aider.
        \end{enumerate}
        Same subject, use ``de + inf'', different subject, use ``que''
        \begin{enumerate}
            \item Je suis content de prendre des vacances.
            \item Il est content que vous veniez demain.
            \item Je suis contente qu'elle soit d'accord avec vous.
            \item Elle a peur de sortir seule.
            \item Nous sommes heureux que vous puissiez nous aider.
        \end{enumerate}

    \section{E7}
        \begin{enumerate}
            \item Elle recevoit une lettre venir de France.
            \item Vous mieux gardez ce sécret.
            \item Il faut que ne accepte pas s'invitation facilement.
            \item << Je vais stationner devant de l'hôtel >> il dit paisiblement à le deux gens covoiturage.
            \item Vous devraiez parler à le professeur de cette événement seule.
        \end{enumerate}
        Correction:
        \begin{enumerate}
            \item Elle reçoit une lettre venant de France.
            \item Vous feriez mieux de garder ce secret.
            \item Il ne faut pas accepter son invitation à la légère.
            \item << Je vais m'arrêter devant l'hôtel >>, dit-il calmement aux deux auto-stoppeuses.
            \item Vous devriez aller parler vous-même de cette affaire au professeur.
            \item Marie Curie est l'une des plus grandes physiciennes du monde. Elle est née en Pologne en 1867. Plus tard, elle est allée faire ses études à Paris, où elle a fait la connaissance de Pierre Curie. Grâce à la découverte du radium, Marie Curie a reçu le prix Nobel de physique en 1903. En 1906, Pierre est mort dans un accident de voiture. Marie a continué leurs recherches sur la radioactivité et est devenue la première femme professeur à la Sorbonne.
        \end{enumerate}

    \section{E8}
        \begin{enumerate}
            \item Ms. Wang n'est pas dans la salle de classe, elle est dans la bibioteque.
            \item Le clé le mette au sous l'étage.
            \item Le stylo est sur la table, vous l'apportez.
            \item Que est entre de le supermarché et du parking?
            \item Le supermarché n'est pas dans le centre-ville, il est dehors de la ville.
            \item Marie n'est pas dernier de la voiture, elle est la devant.
            \item La boulangerie est en face de la poste.
            \item L'industrialisation commencait en milieu de XIX siècle.
        \end{enumerate}
        Correction:
        \begin{enumerate}
            \item Mademoiselle Wang n'est pas dans la salle de classe, elle est à la bibliothèque.
            \item La clé est sous la porte.
            \item Le stylo est sur la table, prends-le !
            \item Qu'est-ce qu'il y a entre le supermarché et le parking ?
            \item Le supermarché n'est pas au centre-ville, il est en dehors de la ville.
            \item Marie n'est pas derrière la voiture, elle est devant la voiture.
            \item La boulangerie est en face de la poste.
            \item L'industrialisation a commencé au milieu du XIXe siècle.
        \end{enumerate}
    
\chapter{Lecon 38}
    \begin{itemize}
        \item être coupé de (separated from)
        \item à quel point (combien) % 多么
        \item quel que (the compond sentence following needs to use subjunctif)
        \item faire partie de (to be part of)
        \item se dérouler (to take place)
        \item ne faire que + inf (just)
        \item suivre l'exemple de (to learn from)
        \item en fonction de (based on, depends on)
        \item se méfier (doubt, distrust)
        \item avant que (before ...) the compond sentence that follow up usually have ``ne'' before the verb.
        \item venir de + inf (to have just done something)
        \item tous ceux qui (all those who)
    \end{itemize}

    \section{E2}
        \begin{enumerate}
            \item Téléphonez-moi avant que vous ne partissiez.
            \item Aidez-la pour qu'elle finisse ce travail.
            \item Lisez lentement afin que nous puissions prendre des notes.
            \item Le temps passe sans qu'on s'en apercevoisse.
            \item Nous ne finirons pas les travaux à moins que vous venissiez.
            \item Je terminerai la traduction pourvu que personne ne venisse me déranger.
            \item Il ne perd pas courage bien qu'il ait des difficultés dans ses études.
            \item Quoique nous ayons peu de temps, nous répondrons à toutes les lettres.
            \item Vous pouvez faire du sport à condition que le médecin le permettiez.
            \item J'accepte de venir à condition que vous invitiez Pierre et Marie.
            \item J'attendrai ici jusqu'à qu'il revenisse.
            \item Le professeur lui explique ce texte jusqu'à ce qu'elle le comprenne.
        \end{enumerate}
        Correction:
        \begin{enumerate}
            \item Téléphonez-moi avant que vous ne partiez.
            \item Aidez-la pour qu'elle finisse ce travail.
            \item Lisez lentement afin que nous puissions prendre des notes.
            \item Le temps passe sans qu'on s'en aperçoive.
            \item Nous ne finirons pas les travaux à moins que vous ne veniez nous aider.
            \item Je terminerai la traduction pourvu que personne ne vienne me déranger.
            \item Il ne perd pas courage bien qu'il ait des difficultés dans ses études.
            \item Quoique nous ayons peu de temps, nous répondrons à toutes les lettres.
            \item Vous pouvez faire du sport à condition que le médecin le permette.
            \item J'accepte de venir à condition que vous invitiez Pierre et Marie.
            \item J'attendrai ici jusqu'à ce qu'il revienne.
            \item Le professeur lui explique ce texte jusqu'à ce qu'elle le comprenne.
        \end{enumerate}

    \section{E3}
        \begin{enumerate}
            \item Rentrons vite avant qu'il ne pleuve.
            \item Vous pouvez partir sous que vous vouliez rester ici.
            \item Ne partez pas à moins que nous le sachions.
            \item Afin que le professeur parle fort, nous ne l'entendons pas bien à cause de ce bruit terrible.
            \item On va acheter des machines modernes bien que la production puisse augmenter.
        \end{enumerate}
        Correction:
        \begin{enumerate}
            \item Rentrons vite avant qu'il ne pleuve.
            \item Vous pouvez partir à moins que vous ne vouliez rester ici.
            \item Ne partez pas sans que nous le sachions.
            \item Bien que le professeur parle fort, nous ne l'entendons pas bien à cause de ce bruit terrible.
            \item On va acheter des machines modernes afin que la production puisse augmenter.
        \end{enumerate}

    \section{E4}
        \begin{enumerate}
            \item Dès que tu auras lu ce roman, tu le passeras à Nicole.
            \item Tu n'es jamais allé au Temple du Ciel, après que tu es à Beijing?
            \item Parce que ils ont fait ce stage, les étuidiants rentrent à l'université.
            \item Pendant que le garagiste réparait notre voiture, nous en avons profité pour visiter la ville à pied.
            \item Ils ne sont pas venus, depuis que ils n'ont pas reçu l'invitation.
        \end{enumerate}
        Correction:
        \begin{enumerate}
            \item Dès que tu auras lu ce roman, tu le passeras à Nicole.
            \item Tu n'es jamais allé au Temple du Ciel, depuis que tu es à Beijing ?
            \item Après qu'ils ont fait ce stage, les étudiants rentrent à l'université.
            \item Pendant que le garagiste réparait notre voiture, nous en avons profité pour visiter la ville à pied.
            \item Ils ne sont pas venus, parce qu'ils n'ont pas reçu l'invitation.
        \end{enumerate}

    \section{E5}
        \begin{enumerate}
            \item Bien que le garçon ne soit pas riche, mais il veut payer pour les deux officiers.
            \item Bien qu'elle ait le temps de faire la cuisine, cependant elle achète des plats tout préparés.
            \item En France, bien qu'on soit méfiant à l'égard du patronat, et pourtant on se rend compte de l'importance des stages professionnels.
            \item Bien que les touristes font vivre ce village, mais les habitants se méfient toujours d'eux.
            \item Bien qu'il ne connaîsse que très peu de mots en fraçais, et pourtant il a l'intention de traduire cet article.
        \end{enumerate}
        Correction:
        \begin{enumerate}
            \item Bien que le garçon ne soit pas riche, il veut payer pour les deux officiers.
            \item Bien qu'elle ait le temps de faire la cuisine, elle achète des plats tout préparés.
            \item En France, bien qu'on soit méfiant à l'égard du patronat, on se rend compte de l'importance des stages professionnels.
            \item Bien que les touristes fassent vivre ce village, les habitants se méfient toujours d'eux.
            \item Bien qu'il ne connaisse que très peu de mots en français, il a l'intention de traduire cet article.
        \end{enumerate}

    \section{E6}
        \begin{enumerate}
            \item Qui que tu sois, tu ne peux pas entrer.
            \item Je ne vous croirai pas, quoi que vous disiez.
            \item Il reste toujours calme, où que soit la situation.
            \item Où que vous habitiez, vous avez besoin d'une voiture.
            \item Quel qu'on dise dans les journaux, je trouve ce film très mauvais.
            \item Quel que soit le moment où vous viendrez, vous serez toujours la bienvenue.
        \end{enumerate}

    \section{E7}
        \begin{enumerate}
            \item Je cherche un endroit où je soie tranquille.
            \item Il y a peu de gens qui sache son adresse.
            \item Il voudrait un appartement qui soit grand et clair.
            \item Pouvez-vous m'indiquer une librairie où on vendiez des livres russes?
        \end{enumerate}
        Correction:
        \begin{enumerate}
            \item Je cherche un endroit où je sois tranquille.
            \item Il y a peu de gens qui sachent son adresse.
            \item Il voudrait un appartement qui soit grand et clair.
            \item Pouvez-vous m'indiquer une librairie où on vende des livres russes ?
        \end{enumerate}

    \section{E8}
        \begin{enumerate}
            \item Avant que vous recontez le résponsable, je vous introduires notre organisation.
            \item Il y ne personne comprend pas cette langue.
            \item Sans que nous invitions, ils attendraient la soirée.
            \item À moins que vous puissiez le voir avant qu'il parti, par allieur le télephoner.
            \item L'enseignent nous explique quelque mots difficilles, afin que nous milieu comprennons cette article.
            \item À condition que le moto ne soit pas trop cher, vous acheterez l'un.
            \item Quoi qu'il essaie, il ne réussire pas.
            \item Qui que vous soyez, vous devraiez trouver l'idée seul.
        \end{enumerate}
        Correction:
        \begin{enumerate}
            \item Avant que vous ne rencontriez le responsable, je vais vous présenter notre entreprise.
            \item Il n'y a personne ici qui comprenne cette langue.
            \item Sans que nous les invitions, ils viendront aussi participer à la soirée.
            \item À moins que vous ne puissiez aller le voir avant son départ, téléphonez-lui.
            \item Le professeur nous a expliqué quelques mots difficiles afin que nous comprenions mieux ce texte.
            \item À condition que la moto ne soit pas trop chère, vous pouvez en acheter une.
            \item Quoi qu'il fasse comme efforts, il ne réussira pas.
            \item Qui que vous soyez, vous devez trouver vous-même une solution.
        \end{enumerate}

        \begin{tabular}{|l|r|r|}
            \hline
            \textbf{Expression} & \textbf{Use subjunctive?} & \textbf{Often uses ne explétif?} \\
            \hline
            \textbf{avant que} & yes & yes \\
            \textbf{à moins que} & yes & yes \\
            \textbf{sans que} & yes & no usually \\
            \textbf{pour que} & yes & no \\
            \textbf{afin que} & yes & no \\
            \textbf{bien que} & yes & no \\
            \textbf{quoique} & yes & no \\
            \textbf{à condition que} & yes & no \\
            \textbf{jusqu'à ce que} & yes & no usually \\
            \hline
        \end{tabular}

\chapter{Lecon 39}
    \begin{enumerate}
        \item attaché à qch (like rely on ...)
        \item estimer que (``think'' for opinion)
        \item estimer (estimate, respected)
        \item être suivi de (after ...)
        \item tomber (meet, for coincidence): La fête nationale tombe un dimanche.
    \end{enumerate}

    \section{E2}
        \begin{enumerate}
            \item Je suis très content qu'elle ait trouvé du travail.
            \item Le professeur a été étonné que nous soyons faits une excellente traduction.
            \item Je n'ai reçu aucunne réponse bien que je lui soie écri il y a trois semaines.
            \item Il est possible que vous soyez vous trompés.
            \item Il a pris des vacances bien qu'il ait eu beaucoup de travail.
            \item C'est la seule personne que nous soyons rencontus.
            \item C'est l'une des plus graves erreurs qu'il soit fait.
            \item Il n'y a personne qui soit visité autant de pays que vous.
            \item C'est la seule résolution qui soit adoptée.
        \end{enumerate}
        Correction:
        \begin{enumerate}
            \item Je suis très content qu'elle ait trouvé du travail.
            \item Le professeur a été étonné que nous ayons fait une excellente traduction.
            \item Je n'ai reçu aucune réponse bien que je lui aie écrit il y a trois semaines.
            \item Il est possible que vous vous soyez trompé(s).
            \item Il a pris des vacances bien qu'il ait eu beaucoup de travail.
            \item C'est la seule personne que nous ayons rencontrée.
            \item C'est l'une des plus graves erreurs qu'il ait faites.
            \item Il n'y a personne qui ait visité autant de pays que vous.
            \item C'est la seule résolution qui ait été adoptée.
        \end{enumerate}

    \section{E3}
        \begin{enumerate}
            \item Il est important que nous nous soyons compris.
            \item Il est dommage qu'il se soit aperçu de son erreur trop tard.
            \item Il est regrettable qui ce discours ait duré trop longtemps.
            \item Il est normal qu'ils ayent abandonnés leur projet.
            \item Il est nécessaire qui leur mariage ait eu lieu dans une église.
        \end{enumerate}
        Correction:
        \begin{enumerate}
            \item Il est important que nous nous soyons compris.
            \item Il est dommage qu'il se soit aperçu de son erreur trop tard.
            \item Il est regrettable que ce discours ait duré trop longtemps.
            \item Il est normal qu'ils aient abandonné leur projet.
            \item Il est nécessaire que leur mariage ait eu lieu dans une église.
        \end{enumerate}

    \section{E4}
        \begin{enumerate}
            \item Non, je ne crois pas que Pierre et Jacques soyent déjà là.
            \item Non, je ne crois pas que le magasin soit encore ouvert.
            \item Non, je ne pense que nous soyons en retard.
            \item Non, je ne crois pas qu'ils aillent souvent au cinéma.
            \item Non, je ne crois pas que nous ayons le temps de visiter cette entreprise.
        \end{enumerate}
        Correction:
        \begin{enumerate}
            \item Non, je ne crois pas que Pierre et Jacques soient déjà là.
            \item Non, je ne crois pas que le magasin soit encore ouvert.
            \item Non, je ne pense pas que nous soyons en retard.
            \item Non, je ne crois pas qu'ils aillent souvent au cinéma.
            \item Non, je ne crois pas que nous ayons le temps de visiter cette entreprise.
        \end{enumerate}

    \section{E5}
        \begin{enumerate}
            \item Je ne crois pas que les étudiants de la classe B aient appris la leçon 40.
            \item Il n'est pas sûr qu'ils aient réussi au concours.
            \item Nous ne pensons pas que Marie ait pu suivre des cours d'anglais.
            \item Je ne pense pas que Paul soit parti.
            \item Il n'est pas certain qu'ils aient été très courageux.
        \end{enumerate}

    \section{E6}
        \begin{enumerate}
            \item Qu'ils l'apprennent, s'ils veulent.
            \item Qu'elle le partie, s'elle le désire.
            \item Qu'il y aille, s'il veut.
            \item Qu'il l'achete, s'il souhaite.
            \item Qu'elle le prenne, s'elle veut.
        \end{enumerate}

    \section{E7}
        \begin{enumerate}
            \item Voilà le village dans lequel j'ai passé mes vacances.
            \item Voilà l'arbre sur lequelle ils ont gravé leur nom.
            \item Voilà les étudiantes avec lequels nous avons voyagé.
            \item Voilà l'avion par lequel elles sont venues.
            \item Voilà les recherches auxquels je m'intéresse beaucoup.
            \item Voilà le dictionnaire sans lequel je ne peux pas travailler.
        \end{enumerate}
        Correction:
        \begin{enumerate}
            \item Voilà le village dans lequel j'ai passé mes vacances.
            \item Voilà l'arbre sur lequel ils ont gravé leur nom.
            \item Voilà les étudiants avec lesquels nous avons voyagé.
            \item Voilà l'avion par lequel elles sont venues.
            \item Voilà les recherches auxquelles je m'intéresse beaucoup.
            \item Voilà le dictionnaire sans lequel je ne peux pas travailler.
        \end{enumerate}

    \section{E8}
        \begin{enumerate}
            \item J'aimerais vous présenter à l'une de mes collègues avec lequel j'ai travaillé pendant deux dans.
            \item Je voudrais vous parler d'un des voyages au cours desquels j'ai pris toutes ces photos.
            \item La police a découvert le nom d'une des personnes chez lequel le bandit s'est caché.
            \item La police a découvrt l'un des clubs dans auquel les terroistes se réunissent souvent.
            \item Je voudrais vous montrer l'une des agences de voyage à lequelle je me suis adressé.
        \end{enumerate}
        Correction:
        \begin{enumerate}
            \item J'aimerais vous présenter l'une de mes collègues avec laquelle j'ai travaillé pendant deux ans.
            \item Je voudrais vous parler d'un des voyages au cours desquels j'ai pris toutes ces photos.
            \item La police a découvert le nom d'une des personnes chez laquelle le bandit s'est caché.
            \item La police a découvert l'un des clubs dans lesquels les terroristes se réunissent souvent.
            \item Je voudrais vous montrer l'une des agences de voyage à laquelle je me suis adressé.
        \end{enumerate}

    \section{E9}
        \begin{enumerate}
            \item Nous sommes invités à le festival, qui aura lieu au mois de juillet. Le fetival à lequel nous sommes invités aura lieu au mois de juillet.
            \item Elle a assisté à les spectacles, qui lui ont beaucoup plu. Les spectacles à lesquels elle a assisté lui ont beaucoup plu.
            \item Nous répondions à les questions, qui ont été posées par des journalistes étrangers. Les questions à lesquelles nous répondions ont été posées par des journalistes étrangers.
            \item Ils ont participé à la manifestation qui a eu lieu hier. La manifesetation à laquel ils ont participé a eu lieu hier.
        \end{enumerate}
        Correction:
        \begin{enumerate}
            \item 
            Nous sommes invités à un festival qui aura lieu au mois de juillet. 
            Le festival auquel nous sommes invités aura lieu au mois de juillet.

            \item 
            Elle a assisté à des spectacles qui lui ont beaucoup plu. 
            Les spectacles auxquels elle a assisté lui ont beaucoup plu.

            \item 
            Nous répondions à des questions qui ont été posées par des journalistes étrangers. 
            Les questions auxquelles nous répondions ont été posées par des journalistes étrangers.

            \item 
            Ils ont participé à une manifestation qui a eu lieu hier. 
            La manifestation à laquelle ils ont participé a eu lieu hier.
        \end{enumerate}

    \section{E10}
        \begin{enumerate}
            \item Il est le plus valuer l'un de présent qui je a reçu.
            \item Qu'il y aller, s'il veut regarder le film.
            \item Je ne crois pas qu'ils trouvent déjà la logement.
            \item Je régrette que tu ne notes pas se nombre de téléphone.
            \item Le festival printemps est un festival traditional le plus reconnu par les chinois.
        \end{enumerate}
        Correction:
        \begin{enumerate}
            \item C'est l'un des cadeaux les plus précieux que j'aie reçus.
            \item Qu'il aille au cinéma, s'il veut voir le film.
            \item Je ne crois pas qu'ils aient déjà trouvé de logement. (use ``de'' for negative sentence)
            \item Je regrette que tu n'aies pas noté son numéro de téléphone. 
            \item La fête du Printemps est la fête traditionnelle préférée des Chinois.
        \end{enumerate}

\chapter{Lecon 40}
    \begin{enumerate}
        \item pauvre + noun (means poor in pity), noun + pauvre (means poor in money)
        \item se passer de (to go without)
        \item par l'intermédiaire de (through / by means of)
        \item protéger qn (qch) de (ou: contre) (to protect ... from / against ...)
        \item ne \dots ni \dots ni \dots (neither nor)
        \item parvenir à (to manage to, to suceed in)
        \item souffrir (suffer)
        \item souffrir qch (to endure something)
        \item souffrir de (to suffer from something)
        \item manquer qch (to miss something)
        \item manquer de qch (to lack something)
        \item manquer à qn (to be missed by someone)
        \item de \dots à \dots (from \dots to \dots)
        \item dès que (as soon as)
    \end{enumerate}

    \section{E2}
        \begin{enumerate}
            \item Il faut que nous préparions les documents pour la conférence.
            \item Il est possible qu'il y ait beaucoup de monde.
            \item Nous espérons que vous arriviez à temps.
            \item Ils demandent que nous partions tout de suite.
            \item Je crois qu'elle prendions l'autobus pour venir.
            \item Je ne crois pas que tu ais de raison.
            \item Il m'a demandé si je avais connaît k'adresse de Mme Dupont.
            \item Je ne sais pas comment aller au Louvre.
            \item Nous souhaitons qu'ils puissent réussir au concours universitaire.
            \item Je serai libre dès que je finisse cette lettre.
            \item Bien que nous restions à 5 semaines à Paris, nous n'avons pas pu tout voir.
            \item Quoi qu'elle dire, elle n'arrivera pas à convaincre son père.
            \item Quoiqu'il pluvue, le médecin sortira pour voir ses malades.
            \item Je ferai avec plaisir ce que vous proposerez, pourvu que vous pussiez me convaincre.
            \item Je vais lire les journaux en attendant que ton père revenisses.
        \end{enumerate}
        Correction:
        \begin{enumerate}
            \item Nous espérons que vous arriverez à temps. (use indicatif, for souhaite que, use subjunctif)
            \item Je crois qu'elle prendra l'autobus pour venir.
            \item Je ne crois pas que tu aies raison. (avoir raison means to be right)
            \item Il m'a demandé si je connaissais l'adresse de Mme Dupont. (demande si is for question, and use imparfait since it is in the past)
            \item Je serai libre dès que j'aurai fini cette lettre. (use indicatif after dès que)
            \item Bien que nous soyons restés cinq semaines à Paris, nous n'avons pas pu tout voir. (subjonctif passé)
            \item Quoi qu'elle dise, elle n'arrivera pas à convaincre son père.
            \item Quoiqu'il pleuve, le médecin sortira pour voir ses malades.
            \item Je ferai avec plaisir ce que vous proposerez, pourvu que vous puissiez me convaincre.
            \item Je vais lire les journaux en attendant que ton père revienne.
        \end{enumerate}

    \section{E3}
        \begin{enumerate}
            \item Que ferez-vous dans deux ans?
            \item J'ai lu ce text dans cinq minutes, et vous?
            \item Il a travaillé au Japon pendant deux ans.
            \item Je n'ai pas vu mes parents pendant deux ans.
            \item On peut aller de New York à Washington en trois heures et demie.
            \item Comment le monde sera-t-il dans 50 ans?
            \item Les étudiants ont fait un stage pendant les vacances.
            \item Tu vas revenir depuis combien de temps?
        \end{enumerate}
        Correction:
        \begin{enumerate}
            \item Que ferez-vous dans deux ans?
            \item J'ai lu ce texte en cinq minutes, et vous?
            \item Il a travaillé au Japon pendant deux ans.
            \item Je n'ai pas vu mes parents depuis deux ans.
            \item On peut aller de New York à Washington en trois heures et demie.
            \item Comment le monde sera-t-il dans 50 ans?
            \item Les étudiants ont fait un stage pendant les vacances.
            \item Tu vas revenir dans combien de temps?
        \end{enumerate}
    
    \section{E4}
        \begin{enumerate}
            \item Excusez-moi de ne pas vous avoir écrit plus tôt.
            \item Il était heureux de avoir fini son travail.
            \item Je suis désolé de vous être dérangés.
            \item Je vous remercie de être viens.
            \item Etes-vous certain de me être laissé la clé?
            \item Jacques est parti sans avoir passée son examen.
        \end{enumerate}
        Correction:
        \begin{enumerate}
            \item Il était heureux d'avoir fini son travail.
            \item Je suis désolé de vous avoir dérangé.
            \item Je vous remercie d'être venu.
            \item Êtes-vous certain de m'avoir laissé la clé ?
            \item Jacques est parti sans avoir passé son examen.
        \end{enumerate}

    \section{E5}
        \begin{enumerate}
            \item Après avoir eu bien joué, les enfants vont se coucher.
            \item Après avoir allé à la banque, elle a fait des achats.
            \item Après avoir assisté à la réunion, nous sommes allés au bar.
            \item Après avoir lu la lettre, elle a pleuré.
            \item Après être arrivé à Paris, il s'est mis à chercher un appartement.
            \item Après avoir lu ce roman, je le lui ai rendu.
            \item Après avoir réfléchi, j'ai pu résoudre mes problèmes.
            \item Après nous avoir été bien amusés, nous sommes rentrés chez nous.
        \end{enumerate}
        Correction:
        \begin{enumerate}
            \item Après avoir bien joué, les enfants vont se coucher.
            \item Après être allée à la banque, elle a fait des achats.
            \item Après nous être bien amusés, nous sommes rentrés chez nous.
        \end{enumerate}

    \section{E6}
        \begin{enumerate}
            \item Après avoir pris son billet, il monte dans le train.
            \item Après l'avoir relué, il envoyé la lettre.
            \item Après aovir réfléchit, elle répond aux questions.
            \item Après avoir hésité pendant deux ans, il choisit une profession.
            \item Après avoir resté six mois à Bruxelles, il écrit un livre sur le Marché Commun.
        \end{enumerate}
        Correction:
        \begin{enumerate}
            \item Après avoir réécrit la lettre, il l'envoie.
        \end{enumerate}



\end{document}